\section{Data and Evaluation}\label{sec:data-eval}

\subsection{Detectors and data}\label{sec:data}
Table~\ref{tab:data} lists the evaluation cells. A cell is one combination of
data, detector, and tracker (the trackers are introduced in
Sect.~\ref{sec:setup}). On MOT17 val-half \cite{milan2016mot16} the trackers
use the published YOLOX-X detections \cite{ge2021yolox}. The emission floor is
the lowest detector score that is passed on; it is 0.01 or 0.1
(Table~\ref{tab:mot17}). On KITTI tracking training \cite{geiger2012kitti},
the trackers run with two detectors that differ from the MOT17 detector:
YOLOv8n \cite{jocher2023yolov8} and RT-DETR-L \cite{zhao2024rtdetr}. To study
the score problem directly, the published MOT17 scores were also changed
before the tracker by four monotone transforms: cubing ($s^3$), halving ($0.5\,s$),
and logit temperatures $T=2$ and $T=0.5$, $\sigma(z/T)$. A temperature
divides the logit $z$ of each score by $T$ before turning it back into a
score: $T=2$ pulls all scores toward 0.5, and $T=0.5$ pushes them toward 0
and 1. The trackers kept their published thresholds. During development,
the layer was also run with ByteTrack on aerial video: the seven sequences of
VisDrone2019-MOT val \cite{zhu2022visdrone} and 40 sequences of its training
set, with YOLOv8n, RT-DETR-L, and (on val) Faster R-CNN detections. These runs
used a class-agnostic protocol and served only to understand the mechanisms;
the same seven val sequences are the evaluation set of our earlier work on
detector-side control (Sect.~\ref{sec:related}). The external systems use the
detections released by their authors: on MOT17 val-half, on the KITTI
tracking sequences of the KITTI MOTS validation split, and on DanceTrack val
\cite{sun2022dancetrack}.

\begin{table*}[t]
\caption{Evaluation cells and their role. Development cells were visible while
the layer was designed. External cells were declared before they were run and
were run once with the frozen configuration: PD-SORT and Hybrid-SORT in the
freeze commit, the others under a protocol registered after the freeze and
before any run.}
\label{tab:data}
\centering\footnotesize
\setlength{\tabcolsep}{3pt}
\begin{tabular}{@{}>{\raggedright\arraybackslash}p{31mm}>{\raggedright\arraybackslash}p{28mm}>{\raggedright\arraybackslash}p{35mm}>{\raggedright\arraybackslash}p{26mm}@{}}
\toprule
Data & Detector & Trackers & Role\\
\midrule
MOT17 val-half (7 sequences) & published YOLOX-X detections & ByteTrack (2 settings), OC-SORT, BoostTrack, SparseTrack & development\\
KITTI tracking training (21 sequences) & RT-DETR-L; YOLOv8n & ByteTrack, BoT-SORT, OC-SORT (each with both detectors) & development (detector change)\\
MOT17 val-half, score transforms & published YOLOX-X, transformed scores & ByteTrack (2 settings), OC-SORT, BoostTrack & development (robustness)\\
VisDrone2019-MOT val (7) and train (40 sequences) & YOLOv8n, RT-DETR-L, Faster R-CNN & ByteTrack & development (aerial video, mechanisms only)\\
MOT17 val-half & released PD-SORT / Hybrid-SORT detections & PD-SORT, Hybrid-SORT & external, declared at the freeze\\
MOT17 val-half; KITTI MOTS val (9 sequences); DanceTrack val (25) & detections released by the authors & C-TWiX & external, declared after the freeze\\
DanceTrack val (25 sequences) & detections released by the authors & TrackTrack & external, declared after the freeze\\
MOT17 val-half & authors' detector & TOPICTrack & external; reproduction FAILED, excluded\\
\bottomrule
\end{tabular}
\end{table*}


\subsection{Evaluation roles and metrics}\label{sec:roles}
Every cell has one of three roles (Table~\ref{tab:data}). Development cells
were visible while \scl{} was designed: MOT17 with four trackers, KITTI with
three trackers and two other detectors, the score transforms, and the VisDrone
runs. Their gains are measured, but they are not held-out evidence. The
external systems were never used during design. Two of them, PD-SORT and
Hybrid-SORT, were chosen and declared in the commit that froze the
configuration. Three more recent systems, C-TWiX, TOPICTrack, and TrackTrack,
were chosen after the freeze under a protocol that was registered before any
run. It fixed the selection criteria, the host contracts, the reproduction
classes, and a prediction of the layer's effect for each system, and it
required every system that met the criteria to be reported, whatever the
result. Each external system was run once with the frozen configuration. For
each external tracker, our reproduction is compared with the authors' number
and marked EXACT (within 0.2 HOTA), CLOSE (within 1.0 HOTA, with a documented
cause that was not tuned), or FAILED; a system with a FAILED reproduction is
excluded from the analysis.

\label{sec:metrics}Table~\ref{tab:metrics} gives the three tracking metrics,
HOTA, multiple object tracking accuracy (MOTA), and identity F1 score (IDF1),
and their definitions. They are built from three error counts: false
positives (FP), missed objects or false negatives (FN), and identity switches
(IDS). They are computed with the TrackEval library (commit
12c8791) and the official ground truth of each benchmark, following
MOTChallenge practice \cite{dendorfer2021motchallenge}. On KITTI tracking
with our two detectors, HOTA is the official value averaged over cars and
pedestrians; for C-TWiX, cars and pedestrians are reported separately, as in
the authors' evaluation.

\begin{table*}[t]
\caption{Tracking metrics used in this paper. GT is the number of
ground-truth boxes; FP, FN, and IDS are the false positives, missed objects
(false negatives), and identity switches, summed over all frames. For IDF1,
each ground-truth track is first matched to at most one output track, and
IDTP, IDFP, and IDFN are the true-positive, false-positive, and
false-negative boxes under this matching. For HOTA, at a localization
threshold $\alpha$, $\mathrm{TP}_\alpha$, $\mathrm{FN}_\alpha$, and
$\mathrm{FP}_\alpha$ are the matched, missed, and unmatched boxes, and
$\mathcal{A}(c)$ is the association accuracy of a matched pair $c$. HOTA,
MOTA, and IDF1 are given in percent.}
\label{tab:metrics}
\centering\footnotesize
\setlength{\tabcolsep}{4pt}
\renewcommand{\arraystretch}{1.6}
\begin{tabular}{@{}l>{\raggedright\arraybackslash}p{68mm}>{\raggedright\arraybackslash}p{40mm}@{}}
\toprule
Metric & Definition & What it measures\\
\midrule
MOTA \cite{bernardin2008clear} & $\displaystyle \mathrm{MOTA}=1-\frac{\mathrm{FN}+\mathrm{FP}+\mathrm{IDS}}{\mathrm{GT}}$ & overall error count; reacts strongly to false positives and missed objects\\
IDF1 \cite{ristani2016idf1} & $\displaystyle \mathrm{IDF1}=\frac{2\,\mathrm{IDTP}}{2\,\mathrm{IDTP}+\mathrm{IDFP}+\mathrm{IDFN}}$ & how long each object keeps the correct identity\\
HOTA \cite{luiten2021hota} & $\displaystyle \mathrm{HOTA}=\frac{1}{19}\sum_{\alpha}\sqrt{\mathrm{DetA}_\alpha\,\mathrm{AssA}_\alpha}$,\newline $\displaystyle \mathrm{DetA}_\alpha=\frac{|\mathrm{TP}_\alpha|}{|\mathrm{TP}_\alpha|+|\mathrm{FN}_\alpha|+|\mathrm{FP}_\alpha|}$,\newline $\displaystyle \mathrm{AssA}_\alpha=\frac{1}{|\mathrm{TP}_\alpha|}\sum_{c\in\mathrm{TP}_\alpha}\mathcal{A}(c)$,\newline with $\alpha=0.05,0.10,\dots,0.95$ & balance of detection (DetA) and association (AssA); the main metric of this paper\\
\bottomrule
\end{tabular}
\end{table*}


All comparisons are paired: each tracker runs with and without \scl{} on
the same sequences. Differences are pooled over sequences, and their 95\%
intervals come from a percentile bootstrap over sequences
\cite{efron1993bootstrap} with 10,000 resamples and seed 42. A difference is
called significant when its interval does not include zero. The intervals
are computed for each comparison separately, without a correction for
multiple comparisons. As a simple
check that does not depend on any distribution, we also count the sequences
where HOTA rose, stayed the same ($|\Delta\mathrm{HOTA}|<0.01$), or fell.

HOTA, MOTA, and IDF1 are given in percent, and higher values are better. FP,
FN, and IDS are counts, and lower values are better. Runtime is given in
milliseconds (ms) per frame, and lower values are better.
